\section{Human review and the semantics of approval}\label{sec:C29:review}
Human review can supply judgment, authority or both. A reviewer may assess whether an artifact is technically sound, while a different principal authorizes publication. Combining these decisions into one button can create ambiguity. The interface should state what is being reviewed, which effect is being authorized, and what evidence remains uncertain. Approval should not be inferred from silence, elapsed time or a historical response to another request.

Review is a limited resource. Long agent-generated explanations can overload the reviewer or hide decisive facts among plausible detail. A useful review surface presents the exact proposed effect, artifact identity, scope change, evidence, failed checks and unresolved consequences. It should distinguish facts derived from execution records from claims generated by the model. This supports informed intervention without pretending that a human inspects every token.

The October Agents SDK release notes explicitly defer automatic compaction while approvals are pending~\cite{executionrelease02}. This is a concrete interaction between context management and authority: changing the context during an unresolved approval can remove or alter the evidence needed to interpret it. The fix is version-specific evidence about one boundary, not a universal solution to approval correctness. Applications must still bind the approval to the action and current policy epoch.

A review experiment can seed defects and vary the presentation of evidence under a fixed time budget. Measure detection, false rejection, decision time, and whether approved actions match what reviewers believed they approved. Also test changes after approval: if the artifact digest or arguments differ at dispatch, the old approval should fail validation. This is a mechanical invariant independent of the reviewer's technical skill.

The check must also have a declared commitment point. Write a grant as $g=(p_g,\mathcal C_g,\alpha_g,t_g)$, binding principal, approved operation envelope, policy epoch, and expiry. For operation $a$ and current principal $p$, a conservative commit-time admission predicate is
\begin{equation}
 \mathbf1\{p=p_g,\ a\in\mathcal C_g\cap F_{\alpha(t)}(t),\
 \alpha(t)=\alpha_g,\ t\leq t_g\}.
\end{equation}
Here an operation includes its arguments, object versions and purpose; $F_{\alpha(t)}(t)$ is the current authoritative feasible set. Requiring epoch equality intentionally invalidates old grants after any governing policy change; a service may instead support a formally specified revalidation rule. Checking this predicate before an unrelated remote write leaves a race between check and use. Enforcement needs a version-conditional or otherwise atomic service contract at the relevant effect boundary. If the service exposes no such primitive, report the residual race and reconciliation behavior rather than claim that an earlier approval proves current authorization. Cancellation cannot retract an effect already committed under the service's own ordering rules.

Escalation policies should distinguish missing preference, missing authority and missing evidence. A user can clarify desired behavior but may not have permission to override an organizational restriction. An administrator can grant a capability without establishing that the proposed action is technically correct. The system should route the question to the appropriate role rather than asking an undifferentiated human oracle to resolve every uncertainty.

